\documentclass[10pt,a4paper]{amsart}
\usepackage{amsmath,amssymb,amsthm,graphicx}
\usepackage{verbatim}
\newenvironment{reference}{\comment}{\endcomment}
%\newenvironment{reference}{}{}
\newenvironment{slogan}{\comment}{\endcomment}
\newenvironment{history}{\comment}{\endcomment}

% For commutative diagrams we use Xy-pic
\usepackage[all]{xy}

% We use 2cell for 2-commutative diagrams.
\xyoption{2cell}
% Theorem environments.
%
\theoremstyle{plain}
\newtheorem{theorem}[subsection]{Theorem}
\newtheorem{proposition}[subsection]{Proposition}
\newtheorem{lemma}[subsection]{Lemma}

\theoremstyle{definition}
\newtheorem{definition}[subsection]{Definition}
\newtheorem{example}[subsection]{Example}
\newtheorem{exercise}[subsection]{Exercise}
\newtheorem{situation}[subsection]{Situation}

\theoremstyle{remark}
\newtheorem{remark}[subsection]{Remark}
\newtheorem{remarks}[subsection]{Remarks}

\numberwithin{equation}{subsection}

% Macros
%
\def\lim{\mathop{\mathrm{lim}}\nolimits}
\def\colim{\mathop{\mathrm{colim}}\nolimits}
\def\Spec{\mathop{\mathrm{Spec}}}
\def\Hom{\mathop{\mathrm{Hom}}\nolimits}
\def\Ext{\mathop{\mathrm{Ext}}\nolimits}
\def\SheafHom{\mathop{\mathcal{H}\!\mathit{om}}\nolimits}
\def\SheafExt{\mathop{\mathcal{E}\!\mathit{xt}}\nolimits}
\def\Sch{\mathit{Sch}}
\def\Mor{\mathop{\mathrm{Mor}}\nolimits}
\def\Ob{\mathop{\mathrm{Ob}}\nolimits}
\def\Sh{\mathop{\mathit{Sh}}\nolimits}
\def\NL{\mathop{N\!L}\nolimits}
\def\CH{\mathop{\mathrm{CH}}\nolimits}
\def\proetale{{pro\text{-}\acute{e}tale}}
\def\etale{{\acute{e}tale}}
\def\QCoh{\mathit{QCoh}}
\def\Ker{\mathop{\mathrm{Ker}}}
\def\Im{\mathop{\mathrm{Im}}}
\def\Coker{\mathop{\mathrm{Coker}}}
\def\Coim{\mathop{\mathrm{Coim}}}

% Boxtimes
%
\DeclareMathSymbol{\boxtimes}{\mathbin}{AMSa}{"02}

%
% Macros for moduli stacks/spaces
%
\def\QCohstack{\mathcal{QC}\!\mathit{oh}}
\def\Cohstack{\mathcal{C}\!\mathit{oh}}
\def\Spacesstack{\mathcal{S}\!\mathit{paces}}
\def\Quotfunctor{\mathrm{Quot}}
\def\Hilbfunctor{\mathrm{Hilb}}
\def\Curvesstack{\mathcal{C}\!\mathit{urves}}
\def\Polarizedstack{\mathcal{P}\!\mathit{olarized}}
\def\Complexesstack{\mathcal{C}\!\mathit{omplexes}}
% \Pic is the operator that assigns to X its picard group, usage \Pic(X)
% \Picardstack_{X/B} denotes the Picard stack of X over B
% \Picardfunctor_{X/B} denotes the Picard functor of X over B
\def\Pic{\mathop{\mathrm{Pic}}\nolimits}
\def\Picardstack{\mathcal{P}\!\mathit{ic}}
\def\Picardfunctor{\mathrm{Pic}}
\def\Deformationcategory{\mathcal{D}\!\mathit{ef}}
\usepackage{cleveref}
\crefname{originallemma}{original lemma}{original lemmas}
\Crefname{originallemma}{Original lemma}{Original lemmas}
\makeatletter
\newlabel{lemma-base-change-module-flat}{{05VW}{source}{}{sourcepointer-1}{}}
\newlabel{lemma-base-change-module-flat@cref}{{[originallemma][0][]05VW}{[1][0][]source}{}{}{}}
\makeatother
\crefname{originallemma}{original lemma}{original lemmas}
\Crefname{originallemma}{Original lemma}{Original lemmas}
\makeatletter
\newlabel{lemma-finite-presentation-local}{{0410}{source}{}{sourcepointer-2}{}}
\newlabel{lemma-finite-presentation-local@cref}{{[originallemma][0][]0410}{[1][0][]source}{}{}{}}
\makeatother
\crefname{originallemma}{original lemma}{original lemmas}
\Crefname{originallemma}{Original lemma}{Original lemmas}
\makeatletter
\newlabel{lemma-flat-local}{{03MM}{source}{}{sourcepointer-3}{}}
\newlabel{lemma-flat-local@cref}{{[originallemma][0][]03MM}{[1][0][]source}{}{}{}}
\makeatother
\crefname{originallemma}{original lemma}{original lemmas}
\Crefname{originallemma}{Original lemma}{Original lemmas}
\makeatletter
\newlabel{lemma-quasi-compact-local}{{03KG}{source}{}{sourcepointer-4}{}}
\newlabel{lemma-quasi-compact-local@cref}{{[originallemma][0][]03KG}{[1][0][]source}{}{}{}}
\makeatother
\crefname{originalproposition}{original proposition}{original propositions}
\Crefname{originalproposition}{Original proposition}{Original propositions}
\makeatletter
\newlabel{morphisms-proposition-generic-flatness-reduced}{{052B}{source}{}{sourcepointer-5}{}}
\newlabel{morphisms-proposition-generic-flatness-reduced@cref}{{[originalproposition][0][]052B}{[1][0][]source}{}{}{}}
\makeatother
\crefname{originallemma}{original lemma}{original lemmas}
\Crefname{originallemma}{Original lemma}{Original lemmas}
\makeatletter
\newlabel{spaces-properties-lemma-etale-image-open}{{06NF}{source}{}{sourcepointer-6}{}}
\newlabel{spaces-properties-lemma-etale-image-open@cref}{{[originallemma][0][]06NF}{[1][0][]source}{}{}{}}
\makeatother
\crefname{originallemma}{original lemma}{original lemmas}
\Crefname{originallemma}{Original lemma}{Original lemmas}
\makeatletter
\newlabel{spaces-properties-lemma-quasi-compact-affine-cover}{{03H6}{source}{}{sourcepointer-7}{}}
\newlabel{spaces-properties-lemma-quasi-compact-affine-cover@cref}{{[originallemma][0][]03H6}{[1][0][]source}{}{}{}}
\makeatother
\crefname{originalsection}{original section}{original sections}
\Crefname{originalsection}{Original section}{Original sections}
\makeatletter
\newlabel{spaces-properties-section-properties-modules}{{05VR}{source}{}{sourcepointer-8}{}}
\newlabel{spaces-properties-section-properties-modules@cref}{{[originalsection][0][]05VR}{[1][0][]source}{}{}{}}
\makeatother
\title{Generic flatness for a finite type morphism of algebraic spaces over a reduced base}
\author{The Stacks Project Authors}
\date{stacks-project:spaces-morphisms.tex:local-snapshot-20261012; GNU FDL 1.2 or later}
\begin{document}
\maketitle
\setcounter{subsection}{1}
\noindent\textit{Editorial note.} All delivered mathematical text is the author's own, copied from the pinned Stacks Project chapter file spaces-morphisms.tex; no statement or proof step was rewritten, reordered or removed.\par\medskip
\noindent\textit{Editorial note.} The proof is the author's reduction to the scheme case; the phrase 'according to our definitions' at the end refers to the definitions of flatness and finite presentation of sheaves on algebraic spaces, which are fixed in the chapter, not to an omitted argument.\par\medskip
\noindent\textit{Editorial note.} The eight references of the counted spans are recorded as non-load-bearing navigation pointers; the argument that consumes them stays inside the counted proof span.\par\medskip
\noindent\textit{Editorial note.} Preamble: the three native excerpt spans of sources/stacks-spaces-morphisms/preamble.tex (lines 8-13, 14-18 and 154-217). The chapter file itself carries no preamble text.\par\medskip
\noindent\textit{Allowed standard prerequisites (editorial).}\begin{itemize}
\item Etale descent of flatness, local finite presentation and quasi-compactness of morphisms. Lemmas lemma-finite-presentation-local, lemma-flat-local and lemma-quasi-compact-local, used at proof lines 6598-6601. Transfers the properties from an etale cover back to X\_W -> W.
\item The image of an etale morphism is open. Properties of Spaces, Lemma spaces-properties-lemma-etale-image-open, used at proof lines 6594-6597. Produces the open subspace W of Y as the image of V'.
\item Flatness and finite presentation of a module after base change by an etale morphism. Lemma lemma-base-change-module-flat and Properties of Spaces, Section spaces-properties-section-properties-modules, used at proof lines 6602-6608. Shows that the descended sheaf is flat and of finite presentation over W.
\item The scheme case of generic flatness. Morphisms, Proposition morphisms-proposition-generic-flatness-reduced, used at proof lines 6633-6638. Is the input that closes the affine case through the etale cover U -> X.
\item Existence of an affine etale cover of a quasi-compact algebraic space. Properties of Spaces, Lemma spaces-properties-lemma-quasi-compact-affine-cover, used at proof lines 6628-6631. Produces the affine scheme U used to reduce to the scheme case.
\end{itemize}\medskip
\ifdefined\hypertarget\hypertarget{sourcepointer-1}{}\fi
\noindent\textit{Source reference.} Original lemma 05VW: focused-seeds/AGP/sources/stacks-snapshot-20261011/tags.txt, source lines 7565--7565. Navigation pointer into the Stacks Project; the standard fact it supplies is applied inside the counted proof span and the argument resting on it stays there.\par\medskip
\ifdefined\hypertarget\hypertarget{sourcepointer-2}{}\fi
\noindent\textit{Source reference.} Original lemma 0410: focused-seeds/AGP/sources/stacks-snapshot-20261011/tags.txt, source lines 5275--5275. Navigation pointer into the Stacks Project; the standard fact it supplies is applied inside the counted proof span and the argument resting on it stays there.\par\medskip
\ifdefined\hypertarget\hypertarget{sourcepointer-3}{}\fi
\noindent\textit{Source reference.} Original lemma 03MM: focused-seeds/AGP/sources/stacks-snapshot-20261011/tags.txt, source lines 4754--4754. Navigation pointer into the Stacks Project; the standard fact it supplies is applied inside the counted proof span and the argument resting on it stays there.\par\medskip
\ifdefined\hypertarget\hypertarget{sourcepointer-4}{}\fi
\noindent\textit{Source reference.} Original lemma 03KG: focused-seeds/AGP/sources/stacks-snapshot-20261011/tags.txt, source lines 4676--4676. Navigation pointer into the Stacks Project; the standard fact it supplies is applied inside the counted proof span and the argument resting on it stays there.\par\medskip
\ifdefined\hypertarget\hypertarget{sourcepointer-5}{}\fi
\noindent\textit{Source reference.} Original proposition 052B: focused-seeds/AGP/sources/stacks-snapshot-20261011/tags.txt, source lines 6564--6564. Navigation pointer into the Stacks Project; the standard fact it supplies is applied inside the counted proof span and the argument resting on it stays there.\par\medskip
\ifdefined\hypertarget\hypertarget{sourcepointer-6}{}\fi
\noindent\textit{Source reference.} Original lemma 06NF: focused-seeds/AGP/sources/stacks-snapshot-20261011/tags.txt, source lines 8534--8534. Navigation pointer into the Stacks Project; the standard fact it supplies is applied inside the counted proof span and the argument resting on it stays there.\par\medskip
\ifdefined\hypertarget\hypertarget{sourcepointer-7}{}\fi
\noindent\textit{Source reference.} Original lemma 03H6: focused-seeds/AGP/sources/stacks-snapshot-20261011/tags.txt, source lines 4555--4555. Navigation pointer into the Stacks Project; the standard fact it supplies is applied inside the counted proof span and the argument resting on it stays there.\par\medskip
\ifdefined\hypertarget\hypertarget{sourcepointer-8}{}\fi
\noindent\textit{Source reference.} Original section 05VR: focused-seeds/AGP/sources/stacks-snapshot-20261011/tags.txt, source lines 7560--7560. Navigation pointer into the Stacks Project; the standard fact it supplies is applied inside the counted proof span and the argument resting on it stays there.\par\medskip
\begin{proposition}
\label{proposition-generic-flatness-reduced}
Let $S$ be a scheme.
Let $f : X \to Y$ be a morphism of algebraic spaces over $S$.
Let $\mathcal{F}$ be a quasi-coherent sheaf of $\mathcal{O}_X$-modules.
Assume
\begin{enumerate}
\item $Y$ is reduced,
\item $f$ is of finite type, and
\item $\mathcal{F}$ is a finite type $\mathcal{O}_X$-module.
\end{enumerate}
Then there exists an open dense subspace $W \subset Y$ such that
the base change $X_W \to W$ of $f$ is flat, locally of finite presentation, and
quasi-compact and such that $\mathcal{F}|_{X_W}$ is flat over $W$ and of
finite presentation over $\mathcal{O}_{X_W}$.
\end{proposition}

\begin{proof}
Let $V$ be a scheme and let $V \to Y$ be a surjective \'etale morphism.
Let $X_V = V \times_Y X$ and let $\mathcal{F}_V$ be the restriction of
$\mathcal{F}$ to $X_V$. Suppose that the result holds for the morphism
$X_V \to V$ and the sheaf $\mathcal{F}_V$. Then there exists an open subscheme
$V' \subset V$ such that $X_{V'} \to V'$ is flat and of finite presentation
and $\mathcal{F}_{V'}$ is an $\mathcal{O}_{X_{V'}}$-module of finite
presentation flat over $V'$. Let $W \subset Y$ be the image of
the \'etale morphism $V' \to Y$, see
Properties of Spaces, Lemma \ref{spaces-properties-lemma-etale-image-open}.
Then $V' \to W$ is a surjective \'etale morphism, hence we see that
$X_W \to W$ is flat, locally of finite presentation, and quasi-compact by
Lemmas \ref{lemma-finite-presentation-local},
\ref{lemma-flat-local}, and
\ref{lemma-quasi-compact-local}.
By the discussion in
Properties of Spaces, Section
\ref{spaces-properties-section-properties-modules}
we see that $\mathcal{F}_W$ is of finite presentation as a
$\mathcal{O}_{X_W}$-module and by
Lemma \ref{lemma-base-change-module-flat}
we see that $\mathcal{F}_W$ is flat over $W$. This argument
reduces the proposition to the case where $Y$ is a scheme.

\medskip\noindent
Suppose we can prove the proposition when $Y$ is an affine scheme.
Let $f : X \to Y$ be a finite type morphism of algebraic spaces
over $S$ with $Y$ a scheme, and let $\mathcal{F}$ be a finite type,
quasi-coherent $\mathcal{O}_X$-module. Choose an affine open covering
$Y = \bigcup V_j$. By assumption we can find dense open $W_j \subset V_j$
such that $X_{W_j} \to W_j$ is flat, locally of finite presentation, and
quasi-compact and such that $\mathcal{F}|_{X_{W_j}}$ is flat over $W_j$
and of finite presentation as an $\mathcal{O}_{X_{W_j}}$-module. In this
situation we simply take $W = \bigcup W_j$ and we win. Hence we reduce
the proposition to the case where $Y$ is an affine scheme.

\medskip\noindent
Let $Y$ be an affine scheme over $S$, let
$f : X \to Y$ be a finite type morphism of algebraic spaces
over $S$, and let $\mathcal{F}$ be a finite type,
quasi-coherent $\mathcal{O}_X$-module. Since $f$ is of finite type
it is quasi-compact, hence $X$ is quasi-compact. Thus we can find
an affine scheme $U$ and a surjective \'etale morphism $U \to X$, see
Properties of Spaces, Lemma
\ref{spaces-properties-lemma-quasi-compact-affine-cover}.
Note that $U \to Y$ is of finite type (this is what it means for
$f$ to be of finite type in this case). Hence we can apply
Morphisms, Proposition \ref{morphisms-proposition-generic-flatness-reduced}
to see that there exists a dense open $W \subset Y$ such that
$U_W \to W$ is flat and of finite presentation and such that
$\mathcal{F}|_{U_W}$ is flat over $W$ and of finite presentation
as an $\mathcal{O}_{U_W}$-module. According to our definitions this means that
the base change $X_W \to W$ of $f$ is flat, locally of finite presentation,
and quasi-compact and $\mathcal{F}|_{X_W}$ is flat over $W$ and of
finite presentation over $\mathcal{O}_{X_W}$.
\end{proof}

\end{document}
